\begin{titlepage}
\thispagestyle{empty}
\setlength{\parindent}{0pt}
\newgeometry{margin=2.4cm}
\begin{tikzpicture}[remember picture,overlay]
  \fill[azulIPN] (current page.north west) rectangle ([yshift=-5.2cm]current page.north east);
  \fill[rojoIPN] ([yshift=-5.2cm]current page.north west) rectangle ([yshift=-5.35cm]current page.north east);
  \node[anchor=north west,text=white,align=left,inner sep=0pt,font=\sffamily] at ([xshift=2.4cm,yshift=-1.5cm]current page.north west)
    {{\Large\bfseries Instituto Politécnico Nacional}\\[5pt]
     {\large Campus Puebla}\\[22pt]
     {\small NASA Space Apps Challenge 2026 \textbullet\ Reto «MODIS and VIIRS Hot Spots»}};
\end{tikzpicture}\par
\vspace*{3.9cm}

{\sffamily\bfseries\fontsize{30}{34}\selectfont\color{azulIPN} Calendario de quemas}\\[10pt]
{\sffamily\Large\color{grisIPN} Armonización de hotspots MODIS y VIIRS y pronóstico para México, 2000--2025}

\vspace{1.0cm}
{\sffamily\bfseries\large\color{rojoIPN} Procedimiento matemático y análisis estadístico}\\[4pt]
{\sffamily\color{grisIPN} Documento técnico para el equipo: cada fórmula se deduce paso a paso, se justifica y se interpreta.}

\vfill
\begin{tcolorbox}[colback=azulclaro,colframe=azulclaro,arc=3pt,left=12pt,right=12pt,top=10pt,bottom=10pt]
{\sffamily\bfseries\color{azulIPN} Equipo}\\[4pt]
\sffamily
Marcos Rodolfo Macías Mondragón\\
Lina Ariel Limón Cortés\\
Joel Morales Gutierrez\\
Alejandro David Zacatelco Escobar
\end{tcolorbox}

\vspace{0.6cm}
{\sffamily\color{grisIPN} Puebla, Puebla \;\textbullet\; noviembre de 2026}
\restoregeometry
\end{titlepage}
